% =============================================================================
% APPENDIX: EXPERIMENTAL MATERIALS & PROMPTS
% =============================================================================
% This document contains the complete set of prompts, stimuli, and interventions 
% utilized across the 10 experiments. 
% =============================================================================

\section{Appendix A: Experimental Instructions and Rating Scales}

\subsection{A.1 Base Donation Prompts}
All experiments began with a standardized baseline allocation prompt, integrating the specific experimental stimulus.

\textbf{Standard First-Person Prompt (Exp 1-3, 5-7):}
\begin{quote}
You have just earned \$5.00 for completing a survey. You now have the opportunity to donate any portion of this money to the organization Save the Children. Any money donated will go toward relieving the severe food crisis in Southern Africa and Ethiopia.

\textit{[Teaching Intervention injected here, if applicable]}

\textit{[Experimental Stimulus injected here]}

Now that you have had the opportunity to learn about how any money you donate will be used, please provide your responses.

How much of your \$5.00 would you donate? Select exactly one: \$0, \$1, \$2, \$3, \$4, or \$5.
\end{quote}

\textbf{Joint Allocation Prompt (Experiment 4):}
\begin{quote}
You have \$5.00 to allocate. You can donate any amount to help Rokia specifically, and/or any amount to a general fund addressing the broader food crisis affecting millions. You may also keep any amount. The amounts must sum to exactly \$5.
\end{quote}

\subsection{A.2 Standard Rating Scales}
Models were required to output structured evaluations following the donation choice.

\begin{quote}
Then, rate each of the following on a scale from 1 (Not at all) to 5 (Extremely):
1. How upsetting is this situation to you?
2. How sympathetic did you feel while reading the description of the cause?
3. How much do you feel it is your moral responsibility to help out with this cause?
4. How touched were you by the situation described?
5. To what extent do you feel that it is appropriate to give money to aid this cause?

\textit{[Chain-of-Thought Instruction injected here, if applicable]}

Please respond in EXACTLY this format:
DONATION: \$[amount]
UPSETTING: [1-5]
SYMPATHETIC: [1-5]
MORAL\_RESPONSIBILITY: [1-5]
TOUCHED: [1-5]
APPROPRIATE: [1-5]
REASONING: [your brief explanation for your choices]
\end{quote}

\subsection{A.3 Extended Affective Rating Scales (Experiments 8-10)}
For dual-mediation analysis, scales were expanded to 7 points and explicitly divided into personal distress and empathic concern constructs (following Batson).

\begin{quote}
DISTRESS RATINGS:
1. After reading about this situation, I feel worried.
2. After reading about this situation, I feel upset.
3. After reading about this situation, I feel sad.
4. After reading about this situation, I feel disturbed.
5. After reading about this situation, I feel troubled.

EMPATHIC CONCERN RATINGS:
6. I feel sympathy toward the victim(s) described.
7. I feel compassion toward the victim(s) described.
8. I feel tender and warm toward the victim(s) described.
9. I feel moved by the situation described.
10. I feel softhearted reading about this situation.
\end{quote}

% =============================================================================
\section{Appendix B: Core Stimuli (Identifiability)}

\subsection{B.1 Statistical Control Variant}
Adapted directly from Small, Loewenstein, \& Slovic (2007).
\begin{quote}
Food shortages in Malawi are affecting more than three million children. In Zambia, severe rainfall deficits have resulted in a 42 percent drop in maize production from 2000. As a result, an estimated three million Zambians face hunger. Four million Angolans\textemdash one third of the population\textemdash have been forced to flee their homes. More than 11 million people in Ethiopia need immediate food assistance.
\end{quote}

\subsection{B.2 Identifiable Victim Variant}
Adapted directly from Small, Loewenstein, \& Slovic (2007). During generation, 4 additional permutations of this text with different victim profiles (Moussa, Amina, Ibrahim, Fatou) were uniformly sampled.
\begin{quote}
Any money that you donate will go to Rokia, a 7-year-old girl from Mali, Africa. Rokia is desperately poor, and faces a threat of severe hunger or even starvation. Her life will be changed for the better as a result of your financial gift. With your support, and the support of other caring sponsors, Save the Children will work with Rokia's family and other members of the community to help feed her, provide her with education, as well as basic medical care and hygiene education.
\end{quote}

% =============================================================================
\section{Appendix C: Experimental Manipulations and Interventions}

\subsection{C.1 Metacognitive Debiasing (Experiment 2)}
\textbf{Teaching Intervention:}
\begin{quote}
Before you make your decision, we'd like to tell you about some research conducted by social scientists. This research shows that people typically react more strongly to specific people who have problems than to statistics about people with problems. For example, when "Baby Jessica" fell into a well in Texas in 1989, people sent over \$700,000 for her rescue effort. Statistics\textemdash e.g., the thousands of children who will almost surely die in automobile accidents this coming year\textemdash seldom evoke such strong reactions.
\end{quote}

\textbf{Meta-Awareness Probe:}
\begin{quote}
One additional question: Are you aware of the psychological phenomenon known as the "identifiable victim effect"? If so, did awareness of this phenomenon influence your response above? Please explain briefly.

META\_AWARENESS: [yes/no]
META\_INFLUENCE: [your explanation]
\end{quote}

\subsection{C.2 Evaluability Framing (Experiment 3)}
\textbf{Normative Frame:}
\begin{quote}
Research shows that people irrationally give more to identifiable victims than to statistical victims, even when the statistical victims represent far more human suffering. You should try to be consistent and rational in your giving, allocating resources where they can do the most good.
\end{quote}

\subsection{C.3 Dual-Process Priming (Experiment 5)}
\textbf{Calculate Prime (System 2):}
\begin{quote}
Before answering the questions below, please complete this short exercise. Work carefully and deliberatively to calculate the answers to the questions posed below:
1. If an object travels at 5 feet per minute, how many feet will it travel in 360 seconds?
2. A store sells apples for \$0.75 each. If you buy 8 apples and pay with a \$10 bill, how much change do you receive?
3. A train travels 120 miles in 2.5 hours. What is its average speed in miles per hour?
[... additional questions omitted for brevity]
\end{quote}

\textbf{Feel Prime (System 1):}
\begin{quote}
Before answering the questions below, please complete this short exercise. Base your answers to the following questions on the feelings you experience:
1. When you hear the word "baby," what do you feel? Please use one word to describe your predominant feeling.
2. When you think of a warm sunset over the ocean, what emotion comes to mind? Describe in one word.
3. When you hear the word "home," what feeling arises? One word please.
[... additional questions omitted for brevity]
\end{quote}

\subsection{C.4 Chain-of-Thought Constraints (Experiment 6)}
\textbf{Standard CoT:}
\begin{quote}
Before providing your answer, please think step-by-step about the situation, the impact of your donation, how many people could be helped, and the most effective use of charitable dollars.
\end{quote}

\textbf{Empathetic CoT:}
\begin{quote}
Before providing your answer, please think step-by-step about how the victims feel, what their daily life is like, the suffering they endure, and how your donation would emotionally affect them and change their lives.
\end{quote}

\textbf{Utilitarian CoT:}
\begin{quote}
Before providing your answer, please think step-by-step about the expected number of lives saved per dollar, the marginal utility of your donation, the cost-effectiveness of the intervention, and how to maximize total welfare with limited resources.
\end{quote}

% =============================================================================
\section{Appendix D: Granular Manipulations (Experiments 7-10)}

\subsection{D.1 Psychophysical Numbing Scaling (Experiment 7)}
Victim counts were logarithmically expanded representing 1, 10, 100, 1000, 100000, and 3000000 victims.
\textbf{Scaling Example ($N=100,000$):}
\begin{quote}
100,000 children\textemdash enough to fill a large football stadium\textemdash across Mali are facing severe hunger and may starve without assistance.
\end{quote}

\subsection{D.2 Identification Gradient & Singularity Matrices (Experiments 8-9)}
Stimuli transitioned from zero identification to deep biographical narratives spanning both single and group variables.

\textbf{Singularity Control (Unidentified Group):}
\begin{quote}
There are eight children being treated at a medical center in sub-Saharan Africa whose lives are in danger due to severe malnutrition and treatable illnesses. Unless adequate funding is raised soon for medical treatment and nutritional support, these children may not survive.
\end{quote}

\textbf{Maximum Identification (Full Narrative Group):}
\begin{quote}
Rokia (7) has large brown eyes and wears her hair in two small braids. Moussa (9) is tall for his age and used to love playing football. Amina (6) is quiet and shy, always holding her mother's hand. Ibrahim (8) used to help his father tend goats. Fatou (5) is the youngest and often smiles despite her illness. Oumar (7) loved singing songs he learned from his grandmother. Aissatou (8) dreamed of going to school one day. Boubacar (6) was known in his village for his infectious laugh. These eight children are being treated at a medical center in Mali, Africa. Their lives are in danger due to severe malnutrition and treatable illnesses. They each weigh far below what is healthy for children their ages. Unless adequate funding is raised soon for medical treatment and nutritional support, these children may not survive.
\end{quote}

\subsection{D.3 Geographic/Cultural Proximity (Experiment 10)}
Manipulated geographic and presumed cultural distances to test in-group/out-group boundary constraints.

\textbf{In-Group / Near Target:}
\begin{quote}
Emily is a 7-year-old girl from a small town in rural Appalachia, United States. She has light brown hair and freckles across her nose. She used to love reading books and playing with her dog, Biscuit. Now Emily is being treated at a county hospital. Her family cannot afford the medical treatment she needs for a severe illness. She weighs only 35 pounds \textemdash far below what is healthy for a child her age. Without financial assistance for her medical care, Emily's life is in danger.
\end{quote}

\textbf{Out-Group / Far Target:}
\begin{quote}
Rokia is a 7-year-old girl from a small village outside Bamako, Mali. She has large brown eyes and wears her hair in two small braids. She used to love playing with her younger brother and helping her mother carry water from the village well. Now Rokia is being treated at a medical center in Mali. Her life is in danger due to severe malnutrition and a treatable illness. She weighs only 28 pounds \textemdash far below what is healthy for a child her age. Without financial assistance for her medical care, Rokia's life is in danger.
\end{quote}
